\section{Buddhism and Tradition}

Much has been made of Guenon's initial rejection of Buddhism as a ``Traditional'' religious form, and his subsequent recantation of that position under the influence of some of his close colleagues. It is important to note the reason for that. Guenon accepted the common belief that Buddhism taught ``anatman'', or the doctrine of ``no-self'', which Guenon rejected as not Traditional. However, he was persuaded that the original and pure Buddhism did indeed teach ``atman'', but was perverted by later Buddhism. Thus, Guenon accepted early Buddhism as Traditional, as did Julius Evola.\footnote{For a rather enthusiastic advocacy of this position, visit \url{https://web.archive.org/web/20100704194342/http://www.attan.com/}}.

However, T R V Murti offers a more subtle interpretation in his masterful study, \textbf{The Central Philosophy of Buddhism}\footnote{\url{https://www.amazon.com/Central-Philosophy-Buddhism-Madhyamika-System/dp/8121510805}}. He points out that the Madhyamika is epistemological and says nothing about the ontological. Thus it denies both the ``doctrine of the self'' as well as the ``doctrine of the no-self''. What is emphasizes in practice is the direct intuition of realtiy as opposed to rational discourse or speculative metaphysics.

But what makes it interesting from our perspective is Murti's suggestion that ``the Madhyamkia Absolutism can serve as the basis for a possible world-culture.'' Murti was well aware of Rene Guenon and his use of the terminology of the Vedanta as a vehicle to express Traditional doctrine. A Hindu himself, Murti recognized that the Vedanta only makes sense within a Hindu framework (this is the reason Guenon stated for his decision to join a Sufi group rather than become a Hindu).

\begin{quotex}
Mahayana absolutism and the Advaita Vedanta are valuable as providing the basis on which a world-culture can be built. It is only absolutism that can make for the fundamental unity of existence and at the same time allow for differences. Catholicity of outlook and tolerance of differences are their very soul; both insist on the universality of the Real and transcendence of the ego-centric standpoint. The Vedanta, however, is traditional in outlook and is bound to the authority of the Veda, and perhaps it presupposes a specific milieu in which alone it can thrive. The Mahayana is quite liberal, and it has proved its capacity to accomodate itself to various religious and social structure, to revitalise and absorb them.
\end{quotex}


This project, if undertaken, would have the advantage of incorporating the techniques and philosophy of a non-dogmatic Buddhism, including Tantrism. Murti does recognize that this won't be the work of a scholar, but rather of a ``spiritual genius''. (This brings to mind Leo Strauss' prediction that the next world class philosopher would arise out of Asia.)


\flrightit{Posted on 2010-03-13 by Cologero}

\begin{center}* * *\end{center}

\begin{footnotesize}
\begin{sffamily}

\texttt{Olavsson on 2015-11-22 at 13:03 said:} 

Discussion continued from here:
\url{http://www.gornahoor.net/?p=8366\&cpage=1\#comment-19710\%5D}

Yes, I certainly think that suggestion is worthy of serious consideration. If Guenon's knowledge about Buddhism had been more profound, he might very well have agreed with Murti's thesis. Now, as we know, what initially had the most formative impact upon Guenon's metaphysical understanding was the Vedanta, and he had been led to the opinion that Buddhism was originally deviant, unorthodox and so on, because of the Buddha Shakyamuni's critical attitude to certain trends within Brahmanism at the time, as well as his differing approach to various questions. As you point out, he was inspired to revise this overly harsh position by Coomaraswamy (himself a Hindu) et al (Marco Pallis?), but in spite of this, Guenon never cared to remedy his lack of in-depth knowledge of the Buddhist forms, which truly constituted one serious whole in his Traditional and metaphysical education. His knowledge of Daoism was much deeper, even though Buddhism was a far more extensive phenomenon of the ``world of Tradition'' in his day, ranging from the exoteric to the esoteric. I'm not saying that an in depth knowledge of all the great traditions is a requirement for true inner realizations of any kind, of course. But since Guenon expressively valued such cross-cultural understanding from a higher point of view, even if eventually devoting oneself to one single path (as he did after being initiated into that Sufi order -- which order was it, by the way?), it is surprising, to say the least, that he wouldn't have been interested in delving more deeply into the Buddhist traditions. As I remember you mentioning in your comment to one of Guenon's letters posted on this site, he couldn't even give a very basic definition of Vajrayana. That is surely no obligatory knowledge for a westerner, but in the case of a man like Guenon, whose interest was Tradition in general and not `only' his own tradition, one should expect better. Julius Evola, of course, was not illiterate in the case of Buddhism, having introduced us to his own analysis of the Buddhist spirit in `The Doctrine of Awakening', and shared some (although the book was more focused on Hindu tantra) insights on Vajrayana in `The Yoga of Power'. With that in view, I am sure Evola would have appreciated such a contribution as Murti's book on madhyamaka.

In order to successfully ``incorporate the techniques and philosophy of a non-dogmatic Buddhism, including Tantrism'', to quote your words, one would probably have to be initiated into such forms and undergo some extensive training. I know that you were initiated by a Tibetan lama in the US. May I ask which lineage(s) he belonged to? If my memory serves me right, you wouldn't really recommend any of us to do the same. Still, to tell the truth, I am wondering whether a stay in India in order to study the Vajrayana tradition might be valuable for my path. I am especially interested in the higher yogic practices, though I am fully aware that this is not exactly entry-level, and would require completing the foundational preparatory practices. However, at this point, the will to devote my life to initiatic realization trumps any other consideration. Since it seems that the Buddhist traditions -- though usually not in any impressive form -- have come to Europe to remain, why let the modernists have a monopoly on it?

This is slightly off-topic, but I came to think of a brief comment in another article where you questioned the authenticity of Evola's perspective on Buddhism. Where do you think he went off the track, so to speak? For me, `The Doctrine of Awakening' is a highly inspiring book, not least because it serves to remind us of the heroic, virile and Aryan qualities that are mostly ignored by modern self-appointed Buddhists. But still, we are dealing with the interpretation of one who was not himself initiated into any Buddhist tradition, which may be biased by his `personal equation' in certain regards. It is not the `gold standard' of authentic Buddhism. For a westerner who wants to understand the spirit of Buddhism, `The Doctrine... ' is a unique book that should be recommended reading for all of a Traditional disposition (not to mention the provoking challenge it would pose to a modernist), but one shouldn't stop there without going to the original sources (as Evola himself did) as well as undertaking a survey of the various Buddhist branches, the living traditions and lineages today, and their historical masters, and also remember that in the case of one specific branch, there are various levels of understanding.

Even when considering the high metaphysics, such as Mahayana's madhyamaka and its emphasis on 'emptiness'/'voidness'/shunyata, one will soon realize that there are varying ways in which this is understood. Do you know of the division between the zhentong and rangtong perspectives on emptiness? I'm not sure if Murti discusses this in his book, as I have yet to read it. A major traditional work on zhentong translated into English from the Tibetan, is available under the title `Mountain Doctrine: Tibet's Fundamental Treatise on Other-Emptiness and the Buddha-Matrix', authored by the Jonang master Dolpopa in the 14th century, numbering some 800 pages. The Jonangpas most strongly emphasized this, as I've understood, but the zhentong view has even been prevalent in orders like Kagyu and Nyingma, the rangtong angle being most pronounced by Gelug (the order of which the Dalai Lama is the most well-known representative). The advantage of zhentong madhyamaka is that it makes it absolutely impossible for any nihilistic misunderstandings to creep in.

\hfill

\end{sffamily}
\end{footnotesize}


